% ---- pg-327 （书页 315） ----

代替参数
\[
\boldsymbol{\mu}=\left[\begin{array}{c}\mu_Y\\ \hdashline \boldsymbol{\mu}_Z\end{array}\right]
\qquad\text{与}\qquad
\boldsymbol{\Sigma}=\left[\begin{array}{c:c}\sigma_{YY}&\boldsymbol{\sigma}_{ZY}'\\ \hdashline \boldsymbol{\sigma}_{ZY}&\boldsymbol{\Sigma}_{ZZ}\end{array}\right]
\]
而得到. \hfill$\blacksquare$

习惯的做法是将均方误差 $\sigma_{YY\cdot Z}=E(Y-\beta_0-\boldsymbol{\beta}'\mathbf{Z})^2$ 的估计量的除数 $n$ 用 $n-(r+1)$ 来代替, 以获得无偏估计量:
\phantomsection\label{eq:7-52}
\begin{equation}
\left(\frac{n-1}{n-r-1}\right)\left(s_{YY}-\mathbf{s}_{ZY}'\mathbf{S}_{ZZ}^{-1}\mathbf{s}_{ZY}\right)
=\frac{\sum\limits_{j=1}^{n}(Y_j-\hat{\beta}_0-\hat{\boldsymbol{\beta}}'\mathbf{Z}_j)^2}{n-r-1}
\tag{7-52}
\end{equation}

\medskip
\noindent\textbf{例 7.12}\quad (单响应回归函数的极大似然估计)

对例 7.6 中给出计算机数据, $Y$(CPU 时间), $Z_1$(命令)和 $Z_2$(增删项目)的 $n=7$ 个观测值所给出的样本均值向量和样本协方差矩阵为
\[
\hat{\boldsymbol{\mu}}=\left[\begin{array}{c}\bar{y}\\ \hdashline \bar{\mathbf{Z}}\end{array}\right]
=\left[\begin{array}{c}150.44\\ \hdashline 130.24\\3.547\end{array}\right]
\]
\[
\mathbf{S}=\left[\begin{array}{c:c}s_{YY}&\mathbf{s}_{ZY}'\\ \hdashline \mathbf{s}_{ZY}&\mathbf{S}_{ZZ}\end{array}\right]
=\left[\begin{array}{c:cc}
467.913&418.763&35.983\\
\hdashline
418.763&377.200&28.034\\
35.983&28.034&13.657
\end{array}\right]
\]
假定 $Y$, $Z_1$ 和 $Z_2$ 是联合正态变量, 求估计的回归函数和估计的均方误差.

结论 7.13 给出极大似然估计值
\[
\hat{\boldsymbol{\beta}}=\mathbf{S}_{ZZ}^{-1}\mathbf{s}_{ZY}
=\begin{bmatrix}0.003\ 128&-0.006\ 422\\-0.006\ 422&0.086\ 404\end{bmatrix}
\begin{bmatrix}418.763\\35.983\end{bmatrix}
=\begin{bmatrix}1.079\\0.420\end{bmatrix}
\]
\[
\hat{\beta}_0=\bar{y}-\hat{\boldsymbol{\beta}}'\bar{\mathbf{z}}
=150.44-[1.079,0.420]\begin{bmatrix}130.24\\3.547\end{bmatrix}
=150.44-142.019
\]
\[
=8.421
\]
于是估计的回归函数为
\[
\hat{\beta}_0+\hat{\boldsymbol{\beta}}'\mathbf{z}=8.42-1.08z_1+0.42z_2
\]
用回归函数预测 $Y$ 的均方误差为
\[
\left(\frac{n-1}{n}\right)\left(s_{YY}-\mathbf{s}_{ZY}'\mathbf{S}_{ZZ}^{-1}\mathbf{s}_{ZY}\right)
\]
\[
=\left(\frac{6}{7}\right)\left(467.913-[418.763,35.983]
\begin{bmatrix}0.003\ 128&-0.006\ 422\\-0.006\ 422&0.086\ 404\end{bmatrix}
\begin{bmatrix}418.763\\35.983\end{bmatrix}\right)
\]
\[
=0.894
\]
\hfill$\blacksquare$

\section{若干个变量的预测}
\phantomsection\label{sec:7.12}

上述结论可直截了当地推广到若干个响应 $Y_1,Y_2,\cdots,Y_m$ 的预测问题. 我们对正态总体给出此种推广.

设
\[
\left[\begin{array}{c}
\underset{(m\times1)}{\mathbf{Y}}\\
\hdashline
\underset{(r\times1)}{\mathbf{Z}}
\end{array}\right]
\quad\text{的联合分布为}\quad N_{m+r}(\boldsymbol{\mu},\boldsymbol{\Sigma})
\]
